\documentclass[10pt,a4paper]{amsart}
\usepackage{amsmath,amssymb,amsthm,graphicx}
% For commutative diagrams we use Xy-pic
\usepackage[all]{xy}

% We use 2cell for 2-commutative diagrams.
\xyoption{2cell}
% Theorem environments.
%
\theoremstyle{plain}
\newtheorem{theorem}[subsection]{Theorem}
\newtheorem{proposition}[subsection]{Proposition}
\newtheorem{lemma}[subsection]{Lemma}

\theoremstyle{definition}
\newtheorem{definition}[subsection]{Definition}
\newtheorem{example}[subsection]{Example}
\newtheorem{exercise}[subsection]{Exercise}
\newtheorem{situation}[subsection]{Situation}

\theoremstyle{remark}
\newtheorem{remark}[subsection]{Remark}
\newtheorem{remarks}[subsection]{Remarks}

\numberwithin{equation}{subsection}

% Macros
%
\def\lim{\mathop{\mathrm{lim}}\nolimits}
\def\colim{\mathop{\mathrm{colim}}\nolimits}
\def\Spec{\mathop{\mathrm{Spec}}}
\def\Hom{\mathop{\mathrm{Hom}}\nolimits}
\def\Ext{\mathop{\mathrm{Ext}}\nolimits}
\def\SheafHom{\mathop{\mathcal{H}\!\mathit{om}}\nolimits}
\def\SheafExt{\mathop{\mathcal{E}\!\mathit{xt}}\nolimits}
\def\Sch{\mathit{Sch}}
\def\Mor{\mathop{\mathrm{Mor}}\nolimits}
\def\Ob{\mathop{\mathrm{Ob}}\nolimits}
\def\Sh{\mathop{\mathit{Sh}}\nolimits}
\def\NL{\mathop{N\!L}\nolimits}
\def\CH{\mathop{\mathrm{CH}}\nolimits}
\def\proetale{{pro\text{-}\acute{e}tale}}
\def\etale{{\acute{e}tale}}
\def\QCoh{\mathit{QCoh}}
\def\Ker{\mathop{\mathrm{Ker}}}
\def\Im{\mathop{\mathrm{Im}}}
\def\Coker{\mathop{\mathrm{Coker}}}
\def\Coim{\mathop{\mathrm{Coim}}}

% Boxtimes
%
\DeclareMathSymbol{\boxtimes}{\mathbin}{AMSa}{"02}

%
% Macros for moduli stacks/spaces
%
\def\QCohstack{\mathcal{QC}\!\mathit{oh}}
\def\Cohstack{\mathcal{C}\!\mathit{oh}}
\def\Spacesstack{\mathcal{S}\!\mathit{paces}}
\def\Quotfunctor{\mathrm{Quot}}
\def\Hilbfunctor{\mathrm{Hilb}}
\def\Curvesstack{\mathcal{C}\!\mathit{urves}}
\def\Polarizedstack{\mathcal{P}\!\mathit{olarized}}
\def\Complexesstack{\mathcal{C}\!\mathit{omplexes}}
% \Pic is the operator that assigns to X its picard group, usage \Pic(X)
% \Picardstack_{X/B} denotes the Picard stack of X over B
% \Picardfunctor_{X/B} denotes the Picard functor of X over B
\def\Pic{\mathop{\mathrm{Pic}}\nolimits}
\def\Picardstack{\mathcal{P}\!\mathit{ic}}
\def\Picardfunctor{\mathrm{Pic}}
\def\Deformationcategory{\mathcal{D}\!\mathit{ef}}
\usepackage{cleveref}
\crefname{originalsection}{original section}{original sections}
\Crefname{originalsection}{Original section}{Original sections}
\makeatletter
\newlabel{sites-modules-section-2-category}{{04IB}{source}{}{sourcepointer-1}{}}
\newlabel{sites-modules-section-2-category@cref}{{[originalsection][0][]04IB}{[1][0][]source}{}{}{}}
\makeatother
\title{Two morphisms of schemes with 2-isomorphic small morphisms of ringed topoi are equal}
\author{The Stacks Project Authors}
\date{stacks-project:etale-cohomology.tex:local-snapshot-20261011; GNU FDL 1.2 or later}
\begin{document}
\maketitle
\noindent\textit{Editorial note.} All delivered mathematical text is the author's own, copied from the pinned Stacks Project chapter file etale-cohomology.tex; no statement or proof step was rewritten, reordered or removed.\par\medskip
\noindent\textit{Editorial note.} The proof is the author's complete argument, including the author's opening remark that careful argument is needed; the Xy-pic diagram is delivered as written.\par\medskip
\noindent\textit{Editorial note.} The label lemma-faithful occurs in no other staged Stacks chapter, so no cross-chapter restatement had to be weighed for this card.\par\medskip
\noindent\textit{Editorial note.} The single labelled reference occurring inside the counted spans points to a chapter of the Project that is not staged in this card's source tree; it is carried in assembly.external\_references as a non-load-bearing navigation pointer recorded at the invoking line.\par\medskip
\noindent\textit{Editorial note.} The preamble is the chapter's own native preamble excerpt already used for this source: the Xy-pic setup (native preamble lines 14-18) and the native theorem declarations with the mathematical macros (native preamble lines 154-217).\par\medskip
\noindent\textit{Allowed standard prerequisites (editorial).}\begin{itemize}
\item Small morphism of ringed topoi attached to a morphism of schemes. Same chapter, the morphism (a\_\{small\}, a\textasciicircum{}sharp) : (X\_etale, O\_X) -> (Y\_etale, O\_Y) associated to a morphism a : X -> Y. Supplies the objects compared by the assumed 2-isomorphism.
\item Subsheaves of the final sheaf on a site and the sheaf h\_V represented by an open subset. Same chapter, the sheaf h\_V on Y\_etale for an open V of Y, a subsheaf of the final sheaf, and the identity a\_\{small\}\textasciicircum{}\{-1\}h\_V = h\_\{a\textasciicircum{}\{-1\}(V)\}. Supplies the argument by which the 2-isomorphism is forced to be the identity on the underlying topological spaces.
\item Sections of the structure sheaf as maps h\_V -> O\_Y, and a\textasciicircum{}sharp as the structural comparison map. Same chapter, the description of a section f in O\_Y(V) by the induced map of sheaves h\_V -> O\_Y and the commutativity of the square defining a 2-isomorphism of ringed topoi. Supplies the diagram whose commutativity gives a\_\{small\}\textasciicircum{}sharp(f) = b\_\{small\}\textasciicircum{}sharp(f).
\end{itemize}\medskip
\ifdefined\hypertarget\hypertarget{sourcepointer-1}{}\fi
\noindent\textit{Source reference.} Original section 04IB: sources/stacks-etale-cohomology/pinned-original-source.tex, source lines 5336--5336. Cross-chapter navigation pointer invoked at proof lines 5335-5336, where the author refers to the definition of a 2-isomorphism of ringed topoi. The target chapter is not staged in this card's source tree, so the pointer records the invoking line; it only resolves the reference for the reader, and no bibliography citation is involved.\par\medskip
\begin{lemma}
\label{lemma-faithful}
Let $X$, $Y$ be schemes.
Any two morphisms $a, b : X \to Y$ of schemes
for which there exists a $2$-isomorphism
$(a_{small}, a_{small}^\sharp) \cong (b_{small}, b_{small}^\sharp)$
in the $2$-category of ringed topoi are equal.
\end{lemma}

\begin{proof}
Let us argue this carefully since it is a bit confusing.
Let $t : a_{small}^{-1} \to b_{small}^{-1}$ be the $2$-isomorphism.
Consider any open $V \subset Y$. Note that $h_V$ is a subsheaf of the
final sheaf $*$. Thus both $a_{small}^{-1}h_V = h_{a^{-1}(V)}$
and $b_{small}^{-1}h_V = h_{b^{-1}(V)}$ are subsheaves of the final sheaf.
Thus the isomorphism
$$
t : a_{small}^{-1}h_V = h_{a^{-1}(V)} \to b_{small}^{-1}h_V = h_{b^{-1}(V)}
$$
has to be the identity, and $a^{-1}(V) = b^{-1}(V)$.
It follows that $a$ and $b$ are equal on underlying topological spaces.
Next, take a section $f \in \mathcal{O}_Y(V)$. This determines and
is determined by a map of sheaves of sets
$f : h_V \to \mathcal{O}_Y$.
Pull this back and apply $t$ to get a commutative diagram
$$
\xymatrix{
h_{b^{-1}(V)} \ar@{=}[r] &
b_{small}^{-1}h_V \ar[d]^{b_{small}^{-1}(f)} & &
a_{small}^{-1}h_V \ar[d]^{a_{small}^{-1}(f)} \ar[ll]^t &
h_{a^{-1}(V)} \ar@{=}[l]
\\
& b_{small}^{-1}\mathcal{O}_Y
\ar[rd]_{b^\sharp}  & &
a_{small}^{-1}\mathcal{O}_Y \ar[ll]^t \ar[ld]^{a^\sharp} \\
& & \mathcal{O}_X
}
$$
where the triangle is commutative by definition of a $2$-isomorphism in
Modules on Sites, Section \ref{sites-modules-section-2-category}.
Above we have seen that the composition of the top horizontal
arrows comes from the identity $a^{-1}(V) = b^{-1}(V)$.
Thus the commutativity of the diagram tells us that
$a_{small}^\sharp(f) = b_{small}^\sharp(f)$ in
$\mathcal{O}_X(a^{-1}(V)) = \mathcal{O}_X(b^{-1}(V))$.
Since this holds for every open $V$ and every $f \in \mathcal{O}_Y(V)$
we conclude that $a = b$ as morphisms of schemes.
\end{proof}

\end{document}
